% !PRINTED 369
\weberpage{369}

\begin{center}
§. 119.

\gesperrt{Convergenz der Näherungsbrüche.}
\end{center}

Wir nehmen jetzt an, dass die Reihe der Zahlen $a_1, a_2, a_3 \ldots$
eine unbegrenzte sei, und bilden die Differenz zweier auf einander
folgender Näherungsbrüche:
\weq{1}{\frac{P_n}{Q_n} - \frac{P_{n-1}}{Q_{n-1}} = \frac{(-1)^n}{Q_n Q_{n-1}}.}

Diese Differenz hat also bei geradem $n$ das positive, bei
ungeradem $n$ das negative Zeichen und nimmt, dem absoluten
Werthe nach, mit unendlich wachsendem $n$ unbegrenzt ab. Bilden
wir noch aus (1):
\weq{2}{\begin{aligned}
\frac{P_n}{Q_n} - \frac{P_{n-2}}{Q_{n-2}} &= \frac{(-1)^n}{Q_n Q_{n-1}} + \frac{(-1)^{n+1}}{Q_{n-1} Q_{n-2}} \\
&= \frac{(-1)^n}{Q_{n-1}}\left(\frac{1}{Q_n} - \frac{1}{Q_{n-2}}\right),
\end{aligned}}

so folgt, da $Q_{n-2} < Q_n$ ist, dass diese Differenz bei geradem $n$
negativ, bei ungeradem $n$ positiv ist.

Hieraus folgt, dass die Reihe der Näherungsbrüche mit geradem
Index
\weq{3}{\frac{P_2}{Q_2}, \quad \frac{P_4}{Q_4}, \quad \frac{P_6}{Q_6} \ldots}

\gesperrt{eine abnehmende}, die Reihe der Näherungsbrüche mit ungeradem
Index
\weq{4}{\frac{P_1}{Q_1}, \quad \frac{P_3}{Q_3}, \quad \frac{P_5}{Q_5} \ldots}

\gesperrt{eine zunehmende ist}.

Nun ist nach \S.\,117
\[
x = \frac{P_n x_n + P_{n-1}}{Q_n x_n + Q_{n-1}},
\]
also
\weq{5}{\frac{P_n}{Q_n} - x = \frac{(-1)^n}{Q_n(Q_n x_n + Q_{n-1})},}

und da $Q_n, Q_{n-1}, x_n$ positiv sind, so sind alle Zahlen der Reihe
(3) grösser als $x$, alle Zahlen der Reihe (4) kleiner als $x$.
